\section{Introduction}
Large Language Models (LLMs)~\citep{radford2018improving,brown2020language,zhang2022opt,openai2023gpt4,touvron2023llama,touvron2023llama2} are becoming ubiquitous, powering many natural language processing applications such as dialog systems~\citep{schulman2022chatgpt,alpaca,vicuna2023}, document summarization~\citep{goyal-durrett-2020-evaluating,zhang2023benchmarking}, code completion~\citep{chen2021evaluating,rozière2023code}  and question answering~\citep{kamalloo2023evaluating}. To unleash the full potential of pretrained LLMs, they should be able to efficiently and accurately perform long sequence generation. For example, an ideal ChatBot assistant can stably work over the content of recent day-long conversations. However, it is very challenging for LLM to generalize to longer sequence lengths than they have been pretrained on, e.g., 4K for Llama-2~\cite{touvron2023llama2}. 

The reason is that LLMs are constrained by the attention window during pre-training. Despite substantial efforts to expand this window size~\citep{chen2023extending,kaiokendev,peng2023yarn} and improve training~\citep{dao2022flashattention,dao2023flashattention2} and inference~\citep{pope2022efficiently,xiao2023smoothquant,anagnostidis2023dynamic,wang2021spatten,zhang2023h2o} efficiency for lengthy inputs, the acceptable sequence length remains intrinsically \emph{finite}, which doesn't allow persistent deployments.

In this paper, we first introduce the concept of LLM streaming applications and ask the question: 
\begin{center}
   \emph{Can we deploy an LLM for infinite-length inputs without sacrificing efficiency and performance?} 
\end{center}
When applying LLMs for infinite input streams, two primary challenges arise:
\vspace{-0.5em}
\begin{enumerate}
    \item During the decoding stage, Transformer-based LLMs cache the Key and Value states (KV) of all previous tokens, as illustrated in Figure~\ref{fig:scheme} (a), which can lead to excessive memory usage and increasing decoding latency~\citep{pope2022efficiently}.
    \item Existing models have limited length extrapolation abilities, i.e., their performance degrades~\citep{alibi,chen2023extending} when the sequence length goes beyond the attention window size set during pre-training.
\end{enumerate}
\vspace{-0.5em}
An intuitive approach, known as window attention~\citep{beltagy2020longformer} (Figure~\ref{fig:scheme} b), maintains only a fixed-size sliding window on the KV states of most recent tokens. Although it ensures constant memory usage and decoding speed after the cache is initially filled, the model collapses once the sequence length exceeds the cache size, \ie, \emph{even just evicting the KV of the first token}, as illustrated in Figure~\ref{fig:nll_curve}.
Another strategy is the sliding window with re-computation (shown in Figure~\ref{fig:scheme} c), which rebuilds the KV states of recent tokens for each generated token. While it offers strong performance, this approach is significantly slower due to the computation of quadratic attention within its window, making this method impractical for real-world streaming applications.
\begin{figure}[t]
    \centering
    \includegraphics[width=\textwidth]{figures/scheme.pdf}
    \vspace{-1.5em}
    \caption{\small
    \textbf{Illustration of \method \vs existing methods.} The language model, pre-trained on texts of length $L$, predicts the $T$th token ($T\gg L$).
    (a) Dense Attention has \(O(T^2)\) time complexity and an increasing cache size. Its performance decreases when the text length exceeds the pre-training text length.
    (b) Window Attention caches the most recent \(L\) tokens' KV. While efficient in inference, performance declines sharply once the starting tokens' keys and values are evicted. (c) Sliding Window with Re-computation rebuilds the KV states from the \(L\) recent tokens for each new token. While it performs well on long texts, its \(O(TL^2)\) complexity, stemming from quadratic attention in context re-computation, makes it considerably slow. (d) \method keeps the \emph{attention sink} (several initial tokens) for stable attention computation, combined with the recent tokens. It's efficient and offers stable performance on extended texts. Perplexities are measured using the Llama-2-13B model on the first book (65K tokens) in the PG-19 test set.}\label{fig:scheme}
    \vspace{-1.5em}
\end{figure}


To understand the failure of window attention, we find an interesting phenomenon of autoregressive LLMs: a surprisingly large amount of attention score is allocated to the initial tokens, irrespective of their relevance to the language modeling task, as visualized in Figure~\ref{fig:attention_visualization}. We term these tokens ``\textbf{attention sinks}". 
Despite their lack of semantic significance, they collect significant attention scores. We attribute the reason to the Softmax operation, which requires attention scores to sum up to one for all contextual tokens. Thus, even when the current query does not have a strong match in many previous tokens, the model still needs to allocate these unneeded attention values somewhere so it sums up to one. The reason behind \textit{initial} tokens as sink tokens is intuitive: initial tokens are visible to almost all subsequent tokens because of the autoregressive language modeling nature, making them more readily trained to serve as attention sinks.

Based on the above insights, we propose \method, a simple and efficient framework that enables LLMs trained with a finite attention window to work on text of infinite length without fine-tuning. \method exploits the fact that attention sinks have high attention values, and preserving them can maintain the attention score distribution close to normal. Therefore, \method simply keeps the attention sink tokens' KV (with just 4 initial tokens sufficing) together with the sliding window's KV to anchor the attention computation and stabilize the model's performance. With \method, models including Llama-2-[7, 13, 70]B, MPT-[7, 30]B, Falcon-[7, 40]B, and Pythia-[2.9,6.9,12]B can reliably model 4 million tokens, and potentially even more. Compared with the only viable baseline, sliding window with recomputation, \method achieves up to 22.2$\times$ speedup, realizing the streaming use of LLMs.

Furthermore, we confirm our attention sink hypothesis and demonstrate that language models can be pre-trained to require only a single attention sink token for streaming deployment.
Specifically, we suggest that an extra learnable token at the beginning of all training samples can serve as a designated attention sink.
By pre-training 160-million parameter language models from scratch, we demonstrate that adding this single sink token preserves the model's performance in streaming cases. This stands in contrast to vanilla models, which necessitate the reintroduction of multiple initial tokens as attention sinks to achieve the same performance level.
\begin{figure}[t]
    \centering
    \includegraphics[width=0.95\textwidth]{figures/attention_weights.pdf}
    \vspace{-1em}
    \caption{\small Visualization of the \emph{average} attention logits in Llama-2-7B over 256 sentences, each with a length of 16. Observations include: (1) The attention maps in the first two layers (layers 0 and 1) exhibit the "local" pattern, with recent tokens receiving more attention.  (2) Beyond the bottom two layers, the model heavily attends to the initial token across all layers and heads.
    }
    \label{fig:attention_visualization}
    \vspace{-2em}
\end{figure}


Finally, we emphasize that \method efficiently generates coherent text from tokens within the KV cache without extending the LLMs' context length. It suits continuous operation needs with minimal memory use and past data reliance. Additionally, \method can complement context extension methods to increase the attendable recent context.
